% ===== Section 3: "Believe in yourself ...", motivation (~18 min) =====
\topicopener{rocket}{Մոտիվացիա}{Հավատա՛ քեզ, ու կկարողանաս ամեն ինչ... ողջամիտ սահմաններում}{Ինչպես են տարբեր մարդիկ գտնում ու օգտագործում մոտիվացիան}{smug}

% Motivational posters
\begin{frame}{Մոտիվացիոն պաստառներ}
  \centering
  \begin{tikzpicture}
    \foreach \x/\e/\q/\ca/\cb in {-3.4/sunrise/{Որքան դժվար է ճանապարհը, այնքան լավ է հասնելը}/peach/lavender,
                                  3.4/mountain/{Ջանքը միայն հիմքն է այն տան, որը դու ես}/mint/butter}{
      \fill[plum, rounded corners=1mm] (\x-2.35,-1.45) rectangle (\x+2.55,1.75);
      \shade[top color=\ca, bottom color=\cb, rounded corners=1mm] (\x-2.45,-1.35) rectangle (\x+2.45,1.85);
      \draw[plum, line width=0.9pt, rounded corners=1mm] (\x-2.45,-1.35) rectangle (\x+2.45,1.85);
      \node at (\x,1.05) {\bigemoji[1.0cm]{\e}};
      \node[text width=4.2cm, align=center, font=\mediumfont\normalsize, text=plum] at (\x,-0.45) {«\q»};}
  \end{tikzpicture}\par\medskip
  \raggedright
  \begin{itemize}
    \item<2-> Մարզասրահում, սրճարանում կամ աշխատավայրի ճաշարանում դրանցից մի քանիսը միշտ կտեսնես
    \item<3-> Զայրույթը կարճ ռեակցիա էր։ Հիմա խոսում ենք \textbf{երկարաժամկետ} մոտիվացիայի մասին, որը ավելի շուտ «մղում» է (drive)
  \end{itemize}
\end{frame}

% What motivation is; procrastination
\begin{frame}{Ի՞նչ է մոտիվացիան}
  \begin{columns}[c]
    \begin{column}{0.52\textwidth}
      \begin{card}[colback=mint!35!cardwhite]
        {\blackfont Մոտիվացիա}\par\smallskip
        Այն «էներգիան», որը պետք է հետաքրքրությունը պահելու ու նախագծի, նպատակի կամ արդյունքի ուղղությամբ աշխատելու համար։
      \end{card}
      \onslide<2->{%
      \begin{card}[colback=brainpink!25!cardwhite]
        {\blackfont Հետաձգում} (procrastination)\par\smallskip
        Մոտիվացիա՝ անելու \textbf{ոչ այն}, ինչ պետք է։ Քանի՜ ռեֆերատ է զոհվել դրան։
      \end{card}}
    \end{column}
    \begin{column}{0.42\textwidth}
      \centering
      \onslide<3->{%
      \bigemoji[1.3cm]{plug}\par\medskip
      \begin{bookquote}
        Գիտեմ՝ ինչի մասին եմ խոսում. այս գիրքը ավարտելու համար ստիպված էի \textbf{անջատել Wi-Fi-ը}։
      \end{bookquote}}
    \end{column}
  \end{columns}
\end{frame}

% Freud's pleasure principle
\begin{frame}{Ֆրեյդի հաճույքի սկզբունքը}
  \begin{columns}[c]
    \begin{column}{0.52\textwidth}
      \begin{itemize}
        \item Վաղ տեսություններից մեկը Զիգմունդ Ֆրեյդինն է. \textbf{հաճույքի սկզբունքը} (pleasure principle)
        \item Կենդանի էակները ձգտում են այն բանին, ինչ \textbf{հաճույք} է տալիս, ու խուսափում են \textbf{ցավից} ու անհարմարությունից
        \item<2-> Որ այդպես լինում է, դժվար է ժխտել։ Կենդանիների ուսուցման փորձերը դա ցույց են տվել
      \end{itemize}
    \end{column}
    \begin{column}{0.42\textwidth}
      \onslide<3->{%
      \begin{bookquote}
        Ֆրեյդը հիմնեց մի ամբողջ ոլորտ, բայց դա չի նշանակում, որ նրա տեսություններն ինքնաբերաբար ճիշտ են։ Ռայթ եղբայրներն էլ ինքնաթիռ հորինեցին, բայց նրանց ինքնաթիռներով Հարավային Ամերիկա չենք թռչում։
      \end{bookquote}}
    \end{column}
  \end{columns}
\end{frame}

% Operant conditioning
\begin{frame}{Առնետը, կոճակը ու պարգևը}
  \centering
  \begin{tikzpicture}
    \node[strip=3.5cm] (a) at (-4.6,0) {\emoji[1.5em]{rat}\\ Առնետը տուփում կոճակ է գտնում ու հետաքրքրությունից սեղմում};
    \node[strip=3.5cm, visible on=<2->] (b) at (0,0) {\emoji[1.5em]{cheese}\\ Սեղմելուց հետո համեղ բան է ընկնում};
    \node[strip=3.5cm, fill=mint!40!cardwhite, visible on=<3->] (c) at (4.6,0) {\emoji[1.5em]{rocket}\\ Հիմա առնետը կոճակը սեղմում է \textbf{անընդհատ}};
    \draw[arr, visible on=<2->] (a) -- (b);
    \draw[arr, visible on=<3->] (b) -- (c);
  \end{tikzpicture}\par\medskip
  \raggedright
  \begin{itemize}
    \item<4-> \textbf{Օպերանտ պայմանավորում} (operant conditioning). պարգևը ուժեղացնում (կամ պատիժը թուլացնում) է իր հետ կապված վարքը
    \item<5-> Երեխային խաղալիք ես խոստանում սենյակը հավաքելու համար, ու նա ավելի պատրաստակամ է նորից հավաքել։ Մեծերի հետ էլ է աշխատում, պարզապես պարգևն է ուրիշը
  \end{itemize}
\end{frame}

% But people do unpleasant things
\begin{frame}{Բայց մարդիկ այդքան պարզ չեն}
  \relax{\small Մենք անընդհատ անում ենք բաներ, որոնք ոչ մի անմիջական հաճույք չեն տալիս}\par\smallskip
  \begin{columns}[T]
    \begin{column}{0.47\textwidth}
      \begin{card}[colback=peach!35!cardwhite]
        \emoji[1.3em]{weightlifter}\ {\blackfont Մարզասրահ}. տանջալի ջանք, իսկ էյֆորիան միշտ չէ, որ գալիս է։ Բըրնեթը դեռ ոչ մի հաճելի փռշտոց էլ չի ստացել այնտեղից
      \end{card}
      \onslide<2->{%
      \begin{card}[colback=mint!35!cardwhite]
        \emoji[1.3em]{gift}\ {\blackfont Բարեգործություն}. սեփական փողը տալիս ես անծանոթների, որոնց երբեք չես տեսնի
      \end{card}}
    \end{column}
    \begin{column}{0.47\textwidth}
      \onslide<3->{%
      \begin{card}[colback=lavender!35!cardwhite]
        \emoji[1.3em]{boss}\ {\blackfont Քծնանք}. շողոքորթում ես տհաճ ղեկավարին՝ առաջխաղացման աղոտ հույսով
      \end{card}}
      \onslide<4->{%
      \begin{card}[colback=butter!45!cardwhite]
        \emoji[1.3em]{books}\ {\blackfont Գրքեր}. կարդում ես գիրք, որ դուր չի գալիս, որովհետև ուզում ես ինչ-որ բան սովորել
      \end{card}}
    \end{column}
  \end{columns}
\end{frame}

% Maslow
\begin{frame}{Մասլոուի պահանջմունքների բուրգը}
  \centering
  {\small 1943. Աբրահամ Մասլոու. հաճույքի փոխարեն՝ \textbf{պահանջմունքներ} (hierarchy of needs)}\par\smallskip
  % Maslow's hierarchy of needs as the usual stepped pyramid (book: "Believe in yourself ...").
% Maslow himself never drew it as a pyramid (see the scienceupdate two frames later). One level per click.
\begin{tikzpicture}[
    lvl/.style={font=\blackfont\footnotesize, text=plum, align=center},
    ex/.style={anchor=west, font=\footnotesize, text=plum, align=left},
  ]
  \def\lvlh{0.84}
  \foreach \i/\c/\name/\e/\examples in {
      0/mint/{Ֆիզիոլոգիական}/burger/{սնունդ, ջուր, օդ},
      1/butter/{Անվտանգություն}/house/{կացարան, անվտանգություն, փող},
      2/peach/{Պատկանելություն}/hug/{հավանություն, աջակցություն, սեր},
      3/lavender/{Հարգանք}/medal/{ուրիշների ու սեփական հարգանքը},
      4/brainpink/{Ինքնաիրացում}/palette/{դառնալ այն, ինչ կարող ես}}{
    \pgfmathsetmacro{\ya}{\i*\lvlh}
    \pgfmathsetmacro{\yb}{(\i+1)*\lvlh}
    \pgfmathsetmacro{\wa}{4.4-\ya*0.68}
    \pgfmathsetmacro{\wb}{4.4-\yb*0.68}
    \pgfmathsetmacro{\ym}{(\ya+\yb)/2}
    \pgfmathsetmacro{\wm}{(\wa+\wb)/2}
    \pgfmathtruncatemacro{\lvlstep}{\i+1}
    \begin{scope}[visible on=<\lvlstep->]
      \filldraw[fill=\c, draw=plum, line width=0.9pt] (-\wa,\ya) -- (\wa,\ya) -- (\wb,\yb) -- (-\wb,\yb) -- cycle;
      \node[lvl] at (0,\ym) {\emoji[1.1em]{\e}\ \name};
      \node[ex] at (\wm+0.25,\ym) {\examples};
    \end{scope}}
\end{tikzpicture}
\par\smallskip
  \onslide<6->{{\small Գաղափարը. նախ ներքևի աստիճանը, հետո հաջորդը, ու այդպես մինչև գագաթ}\par}
\end{frame}

% Self-actualisation + Success Kid
\begin{frame}{Գագաթը. ինքնաիրացում}
  \begin{columns}[c]
    \begin{column}{0.5\textwidth}
      {\small Ինքնաիրացում (self-actualisation). ձգտում՝ դառնալու այն ամենը, ինչ կարող ես}\par\medskip
      \onslide<2->{%
      \begin{bookquote}
        Կարծում ես՝ կարող ես աշխարհի լավագույն նկարիչը լինել։ Ուրեմն մոտիվացված կլինես։ Թեև արվեստը սուբյեկտիվ է, ու գուցե \textbf{արդեն} լավագույնն ես։ Եթե այո, շնորհավորում եմ։
      \end{bookquote}}
    \end{column}
    \begin{column}{0.44\textwidth}
      \centering
      \onslide<3->{%
      \begin{meme}[5.4cm]{success_kid.jpg}
        \memecap[text width=4.6cm]{0.5}{0.9}{Արվեստը սուբյեկտիվ է}
        \memecap[text width=4.6cm]{0.5}{0.1}{Ուրեմն ես արդեն աշխարհի լավագույն նկարիչն եմ}
      \end{meme}}
    \end{column}
  \end{columns}
\end{frame}

% Not so neat + Maslow update
\begin{frame}{Ուղեղը այդքան կոկիկ չէ}
  \begin{columns}[T]
    \begin{column}{0.52\textwidth}
      \begin{itemize}
        \item Մարդիկ վերջին փողը տալիս են անծանոթներին կամ վտանգում իրենց՝ կենդանուն փրկելու համար, թեև կենդանին ոչ հարգանք կտա, ոչ պարգև
        \item<2-> Եթե իշամեղու չէ. սա կխայթի ու չարախնդորեն կծիծաղի \emoji{imp}
        \item<3-> Սեքսը Մասլոուի մոտ ներքևում է, բայց առանց դրա կարելի է ապրել։ Իսկ ինչո՞ւ ենք ուզում. հաճույք, մտերմություն, հարգա՞նք։ Սեքսը \textbf{ամբողջ բուրգով մեկ} է
      \end{itemize}
    \end{column}
    \begin{column}{0.42\textwidth}
      \onslide<4->{%
      \begin{scienceupdate}
        Մասլոուն ինքը \textbf{բուրգ չի նկարել}. այն հայտնվեց 1960-ին՝ կառավարման դասընթացներում։ Իսկ 1943-ին նա գրել էր, որ հերթականությունը «այդքան էլ \textbf{կոշտ} չէ»։
      \end{scienceupdate}}
    \end{column}
  \end{columns}
  \onslide<4->{\source{Maslow, Psychol. Rev., 1943; Bridgman et al., AMLE, 2019}}
\end{frame}

% Intrinsic vs extrinsic
\begin{frame}{Ներքին և արտաքին մոտիվացիա}
  \begin{columns}[T]
    \begin{column}{0.47\textwidth}
      \begin{card}[colback=peach!35!cardwhite]
        {\blackfont Արտաքին} (extrinsic)\par
        \emoji[1.3em]{dollar}\ Վճարում են, որ օգնես տուն տեղափոխվել։ Ձանձրալի է, բայց վճարում են\par
        \onslide<2->{\emoji[1.3em]{cowboy}\ Բոլորը դեղին կովբոյական գլխարկ են կրում, ու դու էլ ես գնում}
      \end{card}
    \end{column}
    \begin{column}{0.47\textwidth}
      \onslide<3->{%
      \begin{card}[colback=mint!35!cardwhite]
        {\blackfont Ներքին} (intrinsic)\par
        \emoji[1.3em]{stethoscope}\ Ինքդ ես որոշել, որ հիվանդներին օգնելն ազնիվ է, ու գնում ես բժշկական\par
        \onslide<4->{Իսկ եթե գնում ես փողի համար, դա ավելի շատ \textbf{արտաքին} է}
      \end{card}}
    \end{column}
  \end{columns}
  \onslide<5->{%
  \begin{chatbox}
    Ինչու՞ ես այսօր այստեղ։ \textbf{Ն}՝ հետաքրքիր է, \textbf{Ա}՝ պետք է, \textbf{ՆԱ}՝ երկուսն էլ
  \end{chatbox}}
\end{frame}

% Self-determination theory
\begin{frame}{Ինչն է մեզ շարժում ներսից}
  \relax{\small Դեսի ու Ռայանի \textbf{ինքնորոշման տեսությունը} (self-determination theory)}\par\smallskip
  \begin{columns}[T]
    \begin{column}{0.31\textwidth}
      \begin{card}[colback=mint!35!cardwhite]
        \emoji[1.4em]{key}\ {\blackfont Ինքնավարություն}\par {\scriptsize autonomy}\par\smallskip
        Ինքդ ես որոշում ու վերահսկում
      \end{card}
    \end{column}
    \begin{column}{0.31\textwidth}
      \begin{card}[colback=butter!45!cardwhite]
        \emoji[1.4em]{trophy}\ {\blackfont Կարողունակություն}\par {\scriptsize competence}\par\smallskip
        Լավ ես անում այն, ինչ անում ես
      \end{card}
    \end{column}
    \begin{column}{0.31\textwidth}
      \begin{card}[colback=lavender!35!cardwhite]
        \emoji[1.4em]{people}\ {\blackfont Կապվածություն}\par {\scriptsize relatedness}\par\smallskip
        Կապ ուրիշների հետ, պատկանելություն*
      \end{card}
    \end{column}
  \end{columns}
  \onslide<2->{%
  \begin{card}[colback=brainpink!25!cardwhite]
    \emoji[1.3em]{eyes}\ Ուսիդ վրայից ամեն քայլդ ուղղող \textbf{միկրոմենեջերը} խլում է երեքն էլ
  \end{card}}
  {\scriptsize * Գրքում՝ «ճանաչում ստանալ արածիդ համար». հեղինակների սահմանումը \textbf{կապն ու պատկանելությունն} է}
\end{frame}

% Overjustification
\begin{frame}{Պարգևը, որ սպանում է հետաքրքրությունը}
  \begin{columns}[c]
    \begin{column}{0.46\textwidth}
      \begin{itemize}
        \item 1973. երեխաները ֆլոմաստերներով նկարում էին։ Ոմանց նախապես \textbf{պարգև էին խոստացել}
        \item<2-> 1-2 շաբաթ անց ավելի շատ նկարում էին \textbf{առանց խոստման} խմբերը
      \end{itemize}
      \onslide<3->{%
      \begin{card}[colback=butter!45!cream]
        \textbf{Գերարդարացման էֆեկտ} (overjustification). խոստացված պարգևը կարող է սպանել ներքին հետաքրքրությունը
      \end{card}}
    \end{column}
    \begin{column}{0.5\textwidth}
      \centering
      \onslide<2->{%
      % Overjustification effect, Lepper, Greene & Nisbett, JPSP 1973 (the book's source): 51 preschoolers drew with
% felt-tip markers; 1-2 weeks later, % of free-play time spent drawing: expected award 8.59, unexpected award 18.09,
% no award 16.73. The book gives no numbers. xtick listed explicitly (see _learnings pgfplots note).
\begin{tikzpicture}
  \begin{axis}[
      width=7.4cm, height=5.0cm,
      ybar, bar width=1.0cm, bar shift=0pt,
      ymin=0, ymax=22,
      symbolic x coords={exp, unexp, none},
      xtick={exp, unexp, none},
      xticklabels={{Պարգև\\ խոստացել էին}, {Անսպասելի\\ պարգև}, {Առանց\\ պարգևի}},
      xticklabel style={align=center, font=\footnotesize, text width=2.2cm},
      ytick={0,5,10,15,20}, yticklabel={\pgfmathprintnumber{\tick}\%},
      tick label style={font=\small, text=plum}, axis line style={plum},
      axis x line*=bottom, axis y line*=left, ymajorgrids, grid style={plum!10},
      nodes near coords={\pgfmathprintnumber[fixed, precision=1]{\pgfplotspointmeta}\%},
      nodes near coords style={font=\blackfont\small, text=plum},
      enlarge x limits=0.25,
    ]
    \addplot[fill=amRed, draw=none] coordinates {(exp,8.59)};
    \addplot[fill=amBlue, draw=none] coordinates {(unexp,18.09)};
    \addplot[fill=amBlue, draw=none] coordinates {(none,16.73)};
  \end{axis}
\end{tikzpicture}
\par
      {\scriptsize Ազատ ժամանակի քանի տոկոսն էին նկարում}}
    \end{column}
  \end{columns}
  \onslide<2->{\source{Lepper, Greene \& Nisbett, JPSP, 1973}}
\end{frame}

% The business world loves it
\begin{frame}{Բիզնեսը շատ ոգևորվեց}
  \begin{columns}[c]
    \begin{column}{0.54\textwidth}
      \begin{itemize}
        \item Եզրակացությունը թվում է ակնհայտ. վարձատրությունը կարող է \textbf{նվազեցնել} մոտիվացիան, իսկ ինքնավարությունը՝ մեծացնել
        \item<2-> Բիզնեսը խանդավառված էր. գիտական հիմնավորում, որ աշխատակցին ավելի լավ է ազատություն տալ, քան \textbf{փող}
        \item<3-> Որոշ հետազոտողներ համաձայն են, բայց հակառակ տվյալներն էլ շատ են
      \end{itemize}
    \end{column}
    \begin{column}{0.42\textwidth}
      \centering
      \onslide<4->{%
      \bigemoji[1.0cm]{moneybag}\par\smallskip
      \begin{bookquote}
        Եթե վճարելը սպանում է մոտիվացիան, միլիոններ ստացող տնօրենները ոչինչ չեն անում։ Ոչ ոք դա չի ասում։ Իսկ միլիարդատերերը կարող են մոտիվացված փաստաբաններ վարձել։
      \end{bookquote}}
    \end{column}
  \end{columns}
\end{frame}

% Self-discrepancy theory
\begin{frame}{Իդեալական «ես»-ը և «պետք է» «ես»-ը}
  \relax{\small Էդվարդ Թորի Հիգինսի տեսությունը (self-discrepancy theory), 1987}\par\smallskip
  \begin{columns}[T]
    \begin{column}{0.31\textwidth}
      \begin{card}
        \emoji[1.4em]{computer}\ {\blackfont Իրական «ես»}\par\smallskip
        Թիկնեղ ծրագրավորող Շոտլանդիայի Ինվերնես քաղաքից
      \end{card}
    \end{column}
    \begin{column}{0.31\textwidth}
      \onslide<2->{%
      \begin{card}[colback=mint!35!cardwhite]
        \emoji[1.4em]{basketball}\ {\blackfont Իդեալական «ես»}\par\smallskip
        Արևայրուք վոլեյբոլիստ Կարիբյան կղզում։ \textbf{Ձգում է} դեպի նպատակը
      \end{card}}
    \end{column}
    \begin{column}{0.31\textwidth}
      \onslide<3->{%
      \begin{card}[colback=peach!35!cardwhite]
        \emoji[1.4em]{salad}\ {\blackfont «Պետք է» «ես»}\par\smallskip
        Յուղոտ չի ուտում, փող չի վատնում, վոլեյբոլ է սովորում ու հետևում Բարբադոսի տների գներին։ \textbf{Հեռու է պահում} սխալներից
      \end{card}}
    \end{column}
  \end{columns}
  \vspace{1mm}
  \onslide<4->{\centering Պիցցա՞։ Դա այն չէ, ինչ «պետք է» անես։ Վերադարձիր աղցաններին \emoji{pizza}\,$\rightarrow$\,\emoji{salad}\par}
\end{frame}

% Locus of control and learned helplessness
\begin{frame}{Ով է ղեկին}
  \begin{columns}[c]
    \begin{column}{0.5\textwidth}
      \begin{itemize}
        \item \textbf{Վերահսկման կենտրոն} (locus of control). որքանով ես զգում, որ դու ես կառավարում իրադարձությունները
        \item Մեկի կարծիքով մոլորակն իր շուրջն է պտտվում, մյուսը միշտ հանգամանքների ողորմածությանն է
        \item<2-> Մասամբ մշակույթից է. արևմտյան մարդը, որին ասել են «կարող ես ունենալ ամեն ինչ», ավելի շատ վերահսկում է զգում, քան տոտալիտար ռեժիմում ապրողը
      \end{itemize}
    \end{column}
    \begin{column}{0.46\textwidth}
      \centering
      \onslide<3->{%
      {\blackfont\small Սովորած անօգնականություն} {\scriptsize (learned helplessness)}\par\smallskip
      \begin{tikzpicture}[font=\footnotesize]
        \node[strip=2.7cm, minimum height=0.9cm] (a) at (0,1.35) {«Ոչինչ չեմ կարող փոխել»};
        \node[strip=2.7cm, minimum height=0.9cm] (b) at (1.8,0) {Չեմ էլ փորձում};
        \node[strip=2.7cm, minimum height=0.9cm] (c) at (0,-1.35) {Ամեն ինչ ավելի է վատանում};
        \node[strip=2.7cm, minimum height=0.9cm] (d) at (-1.8,0) {Ավելի քիչ լավատեսություն};
        \draw[arr] (a) -| (b);
        \draw[arr] (b) |- (c);
        \draw[arr] (c) -| (d);
        \draw[arr] (d) |- (a);
      \end{tikzpicture}\par\smallskip
      {\small Ծանոթ է բոլորին, ով երբևէ \textbf{ծանր բաժանում} է ապրել}}
    \end{column}
  \end{columns}
\end{frame}

% Where motivation lives
\begin{frame}{Որտե՞ղ է մոտիվացիան ուղեղում}
  \begin{columns}[c]
    \begin{column}{0.5\textwidth}
      {\small Հստակ հայտնի չէ, բայց մասնակցում են}\par\smallskip
      \begin{itemize}
        \item միջին ուղեղի \textbf{պարգևի ուղին} (reward pathway)
        \item \textbf{նշաձև մարմինը}. մոտիվացիան զգացմունքային է
        \item \textbf{ճակատային կեղևն} ու այլ «ղեկավար» մասեր. մոտիվացիայի մեծ մասը պլանավորում ու պարգևի սպասում է
      \end{itemize}
    \end{column}
    \begin{column}{0.44\textwidth}
      \onslide<2->{%
      {\small Որոշ գիտնականներ կարծում են, որ կա \textbf{երկու} համակարգ}\par\smallskip
      \begin{card}[colback=lavender!35!cardwhite]
        \emoji[1.3em]{target}\ {\blackfont Բարդ, մտավոր}\par
        կյանքի նպատակներ ու երազանքներ
      \end{card}}
      \onslide<3->{%
      \begin{card}[colback=peach!35!cardwhite]
        \emoji[1.3em]{runner}\ {\blackfont Պարզ, ռեակտիվ}\par
        «Սարսափելի բան, փախի՛ր»։ Կամ՝ «Տե՛ս, տորթ։ Կե՛ր» \emoji{cake}
      \end{card}}
    \end{column}
  \end{columns}
\end{frame}

% Zeigarnik effect
\begin{frame}{Զեյգարնիկի էֆեկտը}
  \begin{columns}[T]
    \begin{column}{0.52\textwidth}
      \begin{itemize}
        \item 1920-ականներ. \textbf{Բլյումա Զեյգարնիկը} նկատեց, որ մատուցողները հիշում են միայն \textbf{չավարտված} պատվերները
        \item<2-> Լաբորատորիայում ընդհատված առաջադրանքներն ավելի լավ էին հիշվում, ու մարդիկ ուզում էին դրանք ավարտել
        \item<3-> Դրա համար են սերիալները կտրվում ամենահետաքրքիր տեղում
      \end{itemize}
    \end{column}
    \begin{column}{0.42\textwidth}
      \onslide<4->{%
      \begin{scienceupdate}
        2025-ի մետավերլուծությունը \textbf{չգտավ}, որ անավարտը ավելի լավ է հիշվում։ Բայց մարդիկ իսկապես հակված են \textbf{վերադառնալու} ընդհատված գործին։
      \end{scienceupdate}}
    \end{column}
  \end{columns}
  \onslide<4->{\source{Ghibellini \& Meier, Humanit. Soc. Sci. Commun., 2025}}
\end{frame}

% Our own cliffhanger
\begin{frame}{Եկեք ստուգենք}
  \begin{columns}[c]
    \begin{column}{0.58\textwidth}
      \begin{card}[colback=butter!45!cream]
        \centering
        {\blackfont\large Ի՞նչ ասաց մեծ ծխնելույզը\\ փոքր ծխնելույզին}\par\smallskip
        \onslide<2->{{\small Պատասխանը՝ մի փոքր ուշ։ Մինչ այդ փորձիր \textbf{չմտածել} դրա մասին}}
      \end{card}
      \vspace{1mm}
      \onslide<3->{%
      \begin{bookquote}
        Մոտիվացնելու երկրորդ լավագույն ձևը գործն անավարտ թողնելն է։ Կա ավելի արդյունավետ ձև, բայց դրա մասին՝ իմ հաջորդ գրքում։
      \end{bookquote}}
    \end{column}
    \begin{column}{0.36\textwidth}
      \centering
      \mascot[2.6cm]{thinking}\par
      \bigemoji[0.9cm]{hourglass}
    \end{column}
  \end{columns}
\end{frame}
